% =====================================================================
% گزارش فنی فارسی - فصل پنجم: توابع سد کنترلی نمایی
% کامپایل با: xelatex -> biber -> xelatex -> xelatex
% =====================================================================
\documentclass[10pt,twocolumn]{article}

% ---------------------- بسته‌های پایه ----------------------
\usepackage{geometry}
\geometry{a4paper, top=2cm, bottom=2.2cm, left=1.8cm, right=1.8cm, columnsep=0.7cm}

\usepackage{xepersian}

% ---------------------- انتخاب فونت (با بازگشت خودکار) ----------------------
% فونت متن اصلی: X Nazanin، بازسازی یونیکدی همان طرح نازنین، هر جا نصب باشد.
% در B Nazanin «ی» و «ک» فارسی با «ي» و «ك» عربی یک گلیف مشترک دارند، پس لایه متن
% PDF (جستجو و کپی) حروف عربی را ثبت می‌کند؛ X Nazanin برای هر کدام و برای همه
% شکل‌های پیوسته‌شان گلیف جداگانه دارد، و «٪» و «٫» را هم دارد که B Nazanin ندارد.
% B Nazanin بازگشت دوم است، برای دستگاهی که X Nazanin ندارد.
\IfFontExistsTF{X Nazanin}
    {\settextfont{X Nazanin}}
    {\IfFontExistsTF{B Nazanin}
        {\settextfont{B Nazanin}}
        {\IfFontExistsTF{XB Niloofar}
            {\settextfont{XB Niloofar}}
            {\settextfont{Vazirmatn}}}}

\IfFontExistsTF{X Nazanin}
    {\setdigitfont{X Nazanin}}
    {\IfFontExistsTF{Geeza Pro}
        {\setdigitfont{Geeza Pro}}
        {\IfFontExistsTF{Damascus}
            {\setdigitfont{Damascus}}
            {\IfFontExistsTF{Vazirmatn}
                {\setdigitfont{Vazirmatn}}
                {\setdigitfont{B Nazanin}}}}}

\IfFontExistsTF{Times New Roman}
    {\setlatintextfont{Times New Roman}}
    {\setlatintextfont{FreeSerif}}

\usepackage{amsmath,amssymb}
\usepackage{graphicx}
\graphicspath{{figures/}}
\usepackage{caption}
\usepackage{subcaption}
\usepackage{booktabs}
\usepackage{float}
\usepackage[hidelinks]{hyperref}
\usepackage{fancyhdr}
\usepackage{titlesec}
\usepackage[backend=biber,style=ieee,sorting=none]{biblatex}
\addbibresource{references.bib}

% ---------------------- دستورهای کمکی ----------------------
% B Nazanin نویسه‌های خط تیره را ندارد؛ آن‌ها را از فونت لاتین می‌گیریم.
\newcommand{\en}[1]{\lr{#1}}
\newcommand{\code}[1]{\lr{\texttt{#1}}}
\newcommand{\md}{\lr{\textemdash}}
\newcommand{\nd}{\lr{\textendash}}

% فهرست مراجع تماماً لاتین است و باید با فونت و جهت لاتین چیده شود.
\AtBeginBibliography{\setlatin}
\renewcommand*{\bibfont}{\latinfont\footnotesize}

% ---------------------- تنظیم عنوان بخش‌ها ----------------------
\titleformat{\section}{\normalfont\bfseries\large}{\thesection.}{0.5em}{}
\titleformat{\subsection}{\normalfont\bfseries}{\thesubsection.}{0.5em}{}

% ---------------------- سربرگ و پابرگ ----------------------
\pagestyle{fancy}
\fancyhf{}
\fancyhead[C]{گزارش فنی \nd{} \small کنترل ایمنی‌محور با توابع سد کنترلی نمایی}
\fancyfoot[C]{\thepage}
\renewcommand{\headrulewidth}{0.4pt}

% ---------------------- اطلاعات سند ----------------------
\newcommand{\reporttitle}{کنترل ایمنی‌محور با توابع سد کنترلی نمایی:\\[0.2em]
بازتولید و اعتبارسنجی عددی فصل پنجم}
\newcommand{\reportauthor}{محمد ملک}
\newcommand{\reportaffil}{پروژه ربات دوپای \en{RABBIT} \nd{} مدل‌سازی، کنترل و بهینه‌سازی مسیر}
\newcommand{\reportemail}{malek.mohammad132@gmail.com}

% =====================================================================
\begin{document}

\onecolumn
\begin{center}
    {\bfseries\Large \reporttitle}\\[0.8em]
    \reportauthor \\
    \small \reportaffil \\
    \small \texttt{\reportemail} \\[0.3em]
    \today
\end{center}
\vspace{1em}

\begin{center}
    {\bfseries\large چکیده}
\end{center}
\noindent
این گزارش فصل پنجم را\md{}توابع سد کنترلی نمایی (\en{ECBF}) برای ایمنی در درجه نسبی
بالا\md{}روی دو سامانه اعتبارسنجی بازتولید می‌کند: زنجیره جرم\nd{}فنر خطی با درجه نسبی $6$ و
پاندول دوحلقه‌ای غیرخطی با عملگر الاستیک و درجه نسبی $4$. سد معکوس تنها در درجه نسبی یک
قید مؤثری بر ورودی می‌سازد؛ روش نمایی با مشتق‌گیری تا مرتبه $r_b$ و جایابی قطب این محدودیت را
برمی‌دارد. در تحقق‌های اسمی، مبنای بدون سد مجموعه امن را ترک می‌کند ($h_{\min}=-0.141$ و
$-0.983$) و هر چهار اجرای \en{ECBF} ایمن می‌مانند. رشته اصلی آن است که تضمینِ زمان پیوسته در
پیاده‌سازی نمونه‌برداری‌شده دوام نمی‌آورد و هر تضمین تنها آنچه را می‌پوشاند که در $h$ نوشته
شده است. روی جرم\nd{}فنر نمونه‌برداری به اندازه $O(dt)$ (شیب $0.998$) از حاشیه می‌کاهد و
ایمنی می‌ماند. روی پاندول، از $36$ اجرای \en{ECBF} با اغتشاش تا $2$ درجه در وضعیت آغازین
تنها $10$ اجرا کامل می‌شود، در حالی که مبنا هر $18$ اجرا را کامل می‌کند؛ از وضعیت آغازین
متحرک، سد در بازوی قائم درجه نسبی خود را از دست می‌دهد؛ و قید ارتفاع با تا کردن بازو تا
$-3.10$ رادیان برآورده می‌شود، خودبرخوردی‌ای که در مدل نوشته نشده است.

\vspace{0.6em}
\noindent\textbf{کلیدواژه‌ها:} تابع سد کنترلی نمایی، ایمنی، تغییرناپذیری پیشرو،
درجه نسبی بالا، تابع لیاپانوف کنترلی، برنامه درجه دوم، جایابی قطب
\vspace{1em}
\twocolumn

% =====================================================================
\section{مقدمه}
% =====================================================================
در طراحی کنترل‌کننده برای سامانه‌های دینامیکی، دو خواسته متفاوت و گاه متعارض
وجود دارد. خواسته نخست \emph{پایداری} است: همگرایی حالت یا خروجی به یک نقطه یا
مسیر مرجع، که ابزار متعارف آن تابع لیاپانوف کنترلی (\en{CLF}) است
\cite{ames2014resclf}. خواسته دوم \emph{ایمنی} است: تضمین اینکه حالت سامانه
هرگز از یک مجموعه مجاز خارج نشود. این دو خواسته یکی نیستند و یکی دیگری را
نتیجه نمی‌دهد؛ یک کنترل‌کننده می‌تواند به‌طور مجانبی پایدار باشد و در گذرا از
مرز مجاز عبور کند.

ابزار متناظر برای خواسته دوم، \textbf{تابع سد کنترلی} (\en{CBF}) است
\cite{ames2014acc,ames2017cbfqp}. جذابیت اصلی هر دو ابزار در یک ویژگی مشترک
ریشه دارد: هر دو به شرطی می‌انجامند که نسبت به ورودی کنترل \emph{اَفین} است و
از این رو می‌تواند به‌صورت یک سطر در یک برنامه درجه دوم (\en{QP}) قرار گیرد
\cite{boyd2004convex}. با این حال، صورت‌بندی متعارف سد\md{}که در این گزارش
«بخش~\en{5.1}» نامیده می‌شود\md{}تنها برای قیدهایی با \emph{درجه نسبی یک} کار می‌کند.

مسئله دقیقاً همین‌جاست. در بسیاری از کاربردهای واقعی، کمیتی که باید محدود شود
از طریق چند انتگرال‌گیر از ورودی کنترل جدا افتاده است: موقعیت مجری نهایی یک
بازوی رباتیک با عملگر الاستیک، یا محل فرود پای آزاد یک ربات دوپا
\cite{nguyen2015footstep}. در چنین حالتی $L_g h \equiv 0$ و سطر سد هیچ اطلاعاتی
درباره ورودی کنترل حمل نمی‌کند.

\textbf{تابع سد کنترلی نمایی} \cite{nguyen2016ecbf} این محدودیت را برطرف
می‌کند: به‌جای مشتق‌گیری یک‌باره از $h$، تا مرتبه $r_b$ (درجه نسبی قید)
مشتق‌گیری می‌شود\md{}یعنی تا جایی که بنا به تعریف، ورودی کنترل نخستین بار ظاهر
می‌شود\md{}و شرط ایمنی روی همان مشتق مرتبه بالا و از طریق جایابی قطب روی یک زنجیره
انتگرال‌گیر اعمال می‌گردد.

هدف این گزارش، \emph{بازتولید عددی} و \emph{اعتبارسنجی} این ساختار است. تمام
مقادیر گزارش‌شده کمیت‌های \textbf{اندازه‌گیری‌شده} از پیاده‌سازی حاضر هستند و
نه مقادیر نقل‌شده از متون.

\textbf{مراجع این فصل.} این گزارش فصل پنجم رساله نگوین \cite{nguyen2017thesis} را بازتولید
می‌کند، که ساختار \en{ECBF} را از \cite{nguyen2016ecbf} می‌آورد. «بخش~\en{5.1}»،
«بخش~\en{5.2}»، «قضیه~\en{5.1}»، «قضیه فرعی~\en{5.2}» و «شکل~\en{5.3}» شماره‌گذاری آن
رساله‌اند و «متن مرجع» همان رساله است.

\textbf{چرا ربات در این فصل نیست.} فصل‌های سوم، چهارم و ششم روی مدلی از \en{RABBIT} کار
می‌کنند که عمداً با تناسبات بدن انسان ساخته شده است ($74$ کیلوگرم، با تنه‌ای $47$ کیلوگرمی)،
زیرا زمینه رساله کنترل برای توان‌بخشی راه‌رفتن با اسکلت خارجی است. فصل پنجم عامداً از آن
ربات استفاده نمی‌کند: روش ابتدا روی دو سامانه‌ای اعتبارسنجی می‌شود که درجه نسبی‌شان معلوم است
و میان خود ادعاهای روش را پوشش می‌دهند، و کاربرد آن بر راه‌رفتن به فصل ششم موکول می‌گردد.
«\emph{ایمن}» در این گزارش یک معنا دارد: $h\ge0$ در همه نمونه‌های اجرا. استخراج گام‌به‌گام
روابط در پیوست~«\ref{app:deriv}» و نام توابع و فایل‌ها در جدول بازتولید
(پیوست~\ref{app:repro5}) آمده است.

% =====================================================================
\section{صورت‌بندی مسئله و مبانی نظری}
% =====================================================================

\subsection{مجموعه امن و تغییرناپذیری پیشرو}
سامانه اَفین نسبت به کنترل را در نظر بگیرید:
\begin{equation}
    \dot{x} = f(x) + g(x)\,u, \qquad y = h_{\mathrm{out}}(x)
    \label{eq:affine}
\end{equation}
با تابع قید $h:\mathcal{X}\to\mathbb{R}$، \textbf{مجموعه امن} چنین تعریف می‌شود:
\begin{equation}
    \mathcal{C} = \{\,x \in \mathcal{X} \;:\; h(x) \ge 0\,\}
    \label{eq:safeset}
\end{equation}
هدف، \textbf{تغییرناپذیری پیشرو} است: اگر $x(0)\in\mathcal{C}$ آنگاه
$x(t)\in\mathcal{C}$ برای هر $t\ge 0$. توجه شود که این خواسته با «ردیابی مرجعی
که اتفاقاً درون مجموعه می‌ماند» یا «جریمه‌کردن خروج از مجموعه» یکی نیست.

در سراسر این پیاده‌سازی یک قرارداد علامت واحد به کار رفته است: قید بالایی
$q\le q_{\max}$ به $h=q_{\max}-q$ و قید پایینی $q\ge q_{\min}$ به
$h=q-q_{\min}$ تبدیل می‌شود، به‌گونه‌ای که مجموعه امن \emph{همواره}
$\{h\ge 0\}$ باشد.

\subsection{سدهای معکوس و محدودیت آن‌ها}
نامزد متعارف سد، تابع معکوس $B=1/h$ است که در مرز به بی‌نهایت میل می‌کند، همراه
با شرط $\dot{B}\le \gamma/B$. با جایگذاری $B=1/h$ و در نتیجه
$\dot{B}=-\dot{h}/h^{2}$، این شرط به صورت بسیار خواناتر زیر در می‌آید که در
پیاده‌سازی نیز از همین فرم استفاده شده است:
\begin{equation}
    \dot{h} \ \ge\ -\gamma\, h^{3},
    \qquad
    \dot{h} = L_f h + L_g h\cdot u
    \label{eq:recip}
\end{equation}
بنابراین این سطر تنها از طریق $L_g h$ بر $u$ اثر می‌گذارد. برای قیدی با درجه
نسبی $r_b>1$، رابطه $L_g h \equiv 0$ برقرار است\md{}این \emph{تعریف} درجه نسبی
است \cite{khalil2002nonlinear}\md{}و سطر سد به گزاره‌ای بدون حضور کنترل تبدیل
می‌شود:
\begin{equation}
    0\cdot u \ \ge\ -\gamma h^{3} - L_f h
    \label{eq:degenerate}
\end{equation}
دو پیامد ممکن است و پیاده‌سازی این دو را از هم تفکیک می‌کند: اگر سمت راست
نامثبت باشد، سطر \emph{به‌طور بدیهی} برقرار است و کنترل‌کننده بی‌سروصدا به یک
\en{CLF-QP} ساده فرو می‌کاهد\md{}شکستی خاموش که تا لحظه عبور از مرز پنهان می‌ماند؛
و اگر سمت راست مثبت شود، سطر \emph{ارضاناپذیر} و برنامه درجه دوم \emph{ناشدنی}
می‌گردد.

محدودیت دومی نیز وجود دارد که کمتر بیان می‌شود: $B=1/h$ تنها روی \emph{درون}
مجموعه، یعنی $\{h>0\}$، یک سد است. در $h\le 0$ مقدار معکوس منفی یا بی‌نهایت است
و هیچ \textbf{رفتار بازیابی} وجود ندارد. در مقابل، شرط نمایی چندجمله‌ای‌ای از
مشتق‌های $h$ است و در $h<0$ نیز کاملاً تعریف‌شده باقی می‌ماند.

\subsection{توابع سد کنترلی نمایی}
راه‌حل، توقف‌نکردن در مشتق نخست است. با مشتق‌گیری تا مرتبه $r_b$ و اعمال تجزیه
خطی‌سازی ورودی-خروجی مجازی:
\begin{align}
    h^{(r_b)} &= L_f^{r_b} h + L_g L_f^{r_b-1} h \cdot u \nonumber\\
              &= L_f^{r_b} h + L_g L_f^{r_b-1} h \cdot (u_{\mathrm{ff}} + A^{-1}\mu)
                 \nonumber\\
              &=: b_0(x) + L_b(x)\,\mu \;=:\; \mu_b
    \label{eq:viol}
\end{align}
کمیت $\mu_b$ نسبت به $\mu$ اَفین است\md{}یک سطر اسکالر، دقیقاً مانند سطر
\en{CLF}\md{}و دینامیک سد به یک زنجیره $r_b$ انتگرال‌گیر در فرم کانونی کنترل‌پذیر
تبدیل می‌شود که در آن \emph{هیچ اثری از سامانه اصلی باقی نمی‌ماند}:
\begin{equation}
    \dot{\eta}_b = F_b \eta_b + G_b \mu_b,\quad
    \eta_b = \big[h,\ \dot{h},\ \ldots,\ h^{(r_b-1)}\big]^{\top}
    \label{eq:chain}
\end{equation}

\subsection{گام از «راندن» به «کران‌گذاری»}
تمام محتوای این ساختار در یک نویسه خلاصه می‌شود. برای \emph{راندن} $h$ به سمت
صفر، قانون کنترل $\mu_b = -K_b\eta_b$ اعمال می‌شود که نتیجه آن
$h(t) = C_b e^{A_b t}\eta_b(0)\to 0$ است. برای \emph{نامنفی نگه‌داشتن} $h$، به
جای آن همان رابطه به صورت یک \emph{الزام} نوشته می‌شود:
\begin{equation}
    \mu_b \ \ge\ -K_b\,\eta_b
    \quad\Longrightarrow\quad
    h(t) \ \ge\ C_b\, e^{A_b t}\,\eta_b(0) \ \ge\ 0
    \label{eq:ecbf}
\end{equation}
تساوی‌ای که یک \emph{قانون کنترل} بود به نامساوی‌ای تبدیل می‌شود که یک
\emph{قید} است، و نمایی‌ای که یک \emph{نرخ همگرایی} بود به یک \textbf{پوش
پایینی} روی سد بدل می‌گردد؛ نام روش از همین‌جا می‌آید. بهره $K_b$ از جایابی قطب
معمولی روی یک زنجیره خطی به دست می‌آید، صرف‌نظر از اینکه سامانه اصلی چه بوده
است. همین ویژگی است که بخش~\en{5.2} را به یک \emph{روش طراحی} تبدیل می‌کند و نه یک
نتیجه وجودی.

\subsection{شرایط قضایا و آنچه واقعاً محدودکننده است}
چهار شرط همراه قضایا هستند. دو شرط ساختاری‌اند و هنگام ساخت بهره اعمال می‌شوند،
و دو شرط به شرط اولیه وابسته‌اند و پیش از هر انتگرال‌گیری در $x_0$ واقعی بررسی
می‌گردند.

نکته اول آنکه \textbf{شرط «تماماً منفی» بودن} $A_b$ (مقادیر ویژه حقیقی و منفی،
و نه صرفاً هورویتس بودن) یک سخت‌گیری تشریفاتی نیست. اثبات قضیه با اعمال
مکرر یک گزاره روی زنجیره مجموعه‌ها انجام می‌شود و هر حلقه از بازگشت
$y_i = \dot{y}_{i-1} + p_i y_{i-1}$ با $p_i$ یک \emph{عدد حقیقی مثبت} استفاده
می‌کند. یک زوج مزدوج مختلط کاملاً هورویتس است و چندجمله‌ای مشخصه درست را نیز
می‌دهد، اما هیچ $p_i$ حقیقی‌ای برای ساخت مجموعه‌های میانی وجود ندارد و استقرا
پله‌ای برای بالا رفتن ندارد. پیاده‌سازی، قطب‌های مختلط را \emph{رد} می‌کند و
مجموعه آزمون بررسی می‌کند که $\mathrm{eig}(F_b - G_b K_b) = -p$ به‌طور دقیق
برقرار باشد و نه صرفاً پایداری حاصل شود. توجه شود که هر دو گزینه
$K_b = \mathrm{fliplr}(\mathrm{poly}(-p))$ و $K_b=\mathrm{poly}(-p)$ هورویتس
هستند و تنها یکی قطب‌های خواسته‌شده را جایابی می‌کند.

نکته دوم آنکه شرط پذیرفتنی بودن قطب‌ها به $x_0$ \emph{وابسته} است. این قاعده را
باید چنین خواند: اگر سد در $x_0$ پیشاپیش به سمت مرز در حرکت باشد، قطب‌ها باید
به‌قدر کافی سریع باشند که آن حرکت را به‌موقع متوقف کرده باشند. در اینجا
قطب‌های کند یک انتخاب محافظه‌کارانه نیستند، بلکه یک انتخاب \emph{ناپذیرفتنی}اند.
این بررسی پیش از انتگرال‌گیری انجام می‌شود، زیرا اجرا از یک $x_0$ ناپذیرفتنی
کنترل‌کننده‌ای تولید می‌کند که سطر خود را در تک‌تک لحظه‌ها وفادارانه اعمال
می‌کند و با این حال اجازه می‌دهد $h$ منفی شود\md{}که دقیقاً شبیه یک اشکال در سد به
نظر می‌رسد و چنین نیست.

\textbf{حاشیه این شرط در اجراهای این گزارش.} هر چهار پیکربندی
(با بررسی شرط روی $x_0$ خودِ هر اجرا) شرط را برقرار می‌کنند، و
\emph{با حاشیه کامل}: در هر چهار مورد $x_0$ از سکون آغاز می‌شود، پس $\dot h(0)=0$
و کران پایین هر قطب صفر است. بنابراین حاشیه هر قطب برابر خودِ آن قطب است
(جدول~\ref{tab:adm}). به بیان دیگر این اجراها محدودیت وابستگی به شرط اولیه را
\emph{نمی‌آزمایند}؛ آن محدودیت وقتی خود را نشان می‌دهد که سد در لحظه آغاز به
سمت مرز حرکت کند. پله نخست زنجیره سرعت بسته‌شدن را به $-\dot h(0)\le p_1h(0)$
محدود می‌کند\md{}$3.78$ و $3.60$ متر بر ثانیه برای جرم\nd{}فنر\md{}اما این تنها یک
سقف است و پله‌های بالاتر بسیار زودتر می‌بندند. اگر هر سه ارابه با سرعت یکسان $s$ به
سمت مرز حرکت کنند، فنرها کشیده نمی‌شوند و همه مشتق‌های بالاتر $h$ در $x_0$ صفرند، پس
زنجیره شکل بسته دارد: $y_i(x_0)\ge0$ اگر و تنها اگر $s\le h(0)/\sum_{j\le i}1/p_j$.
پله‌های یکم تا پنجم، که دقیقاً بررسی‌پذیرند، تنها $0.894$ و $0.852$ متر بر ثانیه را
می‌پذیرند، حدود یک‌چهارم سقف پله نخست. در آونگ، $x_0$ بالاترین نقطه دایره مجری
نهایی است، پس برای هر چرخشی $\dot h(0)=0$ و پله نخست هرگز فعال نمی‌شود؛ زنجیره در
آنجا تنها از راه جمله‌های مرکزگرا در پله‌های بالاتر می‌شکند. بخش~\ref{sec:moving} هر
دو سامانه را از $x_0$ متحرک اجرا می‌کند.

\begin{table}[H]
    \centering
    \caption{شرط پذیرفتنی بودن قطب‌ها در $x_0$ هر اجرا: مقدار سد، تنگ‌ترین قطب،
    کران پایین لازم و حاشیه}
    \label{tab:adm}
    \small
    \begin{tabular}{lccccc}
        \toprule
        اجرا & $r_b$ & $h(x_0)$ & $p_1$ & کران لازم & حاشیه \\
        \midrule
        جرم\nd{}فنر، $x_3^{\max}=3.15$ & $6$ & $3.15$ & $1.20$ & $0$ & $1.20$ \\
        جرم\nd{}فنر، $x_3^{\max}=3.00$ & $6$ & $3.00$ & $1.20$ & $0$ & $1.20$ \\
        آونگ، $p_{2\min}=-1.0$ & $4$ & $3.00$ & $5.00$ & $0$ & $5.00$ \\
        آونگ، $p_{2\min}=-0.5$ & $4$ & $2.50$ & $5.00$ & $0$ & $5.00$ \\
        \bottomrule
    \end{tabular}
\end{table}

در هر چهار مورد همه $y_i(x_0)$ نامنفی‌اند (برای جرم\nd{}فنر از $3.15$ تا $33.9$ و
برای آونگ از $3$ تا $4950$)، پس $x_0$ در همه مجموعه‌های زنجیره قرار دارد و شرطی که
قضیه~$5.1$ بر $x_0$ می‌گذارد بدون استثنا برقرار است. فرض درجه نسبی آن در آونگ داستان
دیگری است (بخش~\ref{sec:moving}).

\subsection{برنامه درجه دوم ایمنی-پایداری}
کنترل‌کننده نهایی از حل برنامه زیر در هر لحظه نمونه‌برداری به دست می‌آید:
\begin{equation}
\begin{aligned}
    \min_{\mu,\,\delta}\quad & \mu^{\top}\mu + p_{\delta}\,\delta^{2}\\
    \text{با قید}\quad
        & \dot{V}(\eta,\mu) + \lambda V(\eta) \le \delta && \text{(پایداری)}\\
        & A_c(x)\,\mu \le b_c(x) && \text{(قید ورودی)}\\
        & \mu_b \ge -K_b\,\eta_b && \text{(ایمنی)}\\
        & h^{(r_b)}(x,\mu) = \mu_b &&
\end{aligned}
\label{eq:qp}
\end{equation}

\textbf{سطر سد هرگز شل نمی‌شود.} تنها سطر \en{CLF} متغیر کمکی $\delta$ (اسلک)
دریافت می‌کند. هرجا این دو با هم در تعارض باشند\md{}و روی پاندول واقعاً در تعارض
هستند\md{}ایمنی برنده می‌شود و ردیابی به تعویق می‌افتد. قید ایمنی‌ای که حاضر باشید
برای هموارتر شدن ورودی نقض کنید، اصلاً قید ایمنی نیست. همین عدم‌تقارن است که
دلیل وجود دو سطر مجزا در~\eqref{eq:qp} به‌جای یک تابع هدف وزن‌دار است.

% =====================================================================
\section{سامانه‌های اعتبارسنجی}
% =====================================================================
دو سامانه به‌گونه‌ای انتخاب شده‌اند که میان خود، ادعاهای روش را پوشش دهند: یکی
خطی و دیگری غیرخطی، یکی تک‌ورودی و دیگری چندورودی، در دو درجه نسبی متفاوت.
مشخصات کامل در جدول~\ref{tab:plants} آمده است.

نکته ساختاری تعیین‌کننده آن است که در \emph{هر دو} سامانه، \textbf{کمیت مقیدشده
همان کمیت کنترل‌شده است}. بنابراین قید ایمنی، درجه نسبی کامل سامانه را به ارث
می‌برد\md{}در آونگ جز جایی که $\nabla p_2^y=0$ است، یعنی بازوی قائم، که وضعیت آغاز را هم
دربرمی‌گیرد (بخش~\ref{sec:moving})\md{}و هیچ بازنویسی ارزانی وجود ندارد که در آن سد به
درجه نسبی یک برسد. این
دقیقاً همان وضعیتی است که سد معکوس بخش~\en{5.1} توان پرداختن به آن را ندارد.

\begin{table}[H]
    \centering
    \caption{مشخصات دو سامانه اعتبارسنجی و ثوابت طراحی}
    \label{tab:plants}
    \small
    \begin{tabular}{lcc}
        \toprule
        ویژگی & جرم-فنر سری & پاندول دوحلقه‌ای \\
        \midrule
        دینامیک            & خطی      & غیرخطی \\
        $n_x / n_u / n_y$  & $6/1/1$  & $8/2/2$ \\
        \textbf{درجه نسبی} & $\mathbf{6}$ & $\mathbf{4}$ \\
        قید ایمنی          & $x_3 \le x_3^{\max}$ & $p_2^{y} \ge p_{2\min}$ \\
        دینامیک صفر        & ندارد    & ندارد \\
        مأموریت            & $x_3: 0\to 3$ & $\theta_1: -\pi\to\pi$ \\
        قطب‌های \en{ECBF}  & $0.12\!\cdot\![10..15]$ & $[5,\,5.5,\,6,\,10]$ \\
        $\lambda$ (نرخ \en{CLF}) & $0.20$ & $1.0$ \\
        افق شبیه‌سازی      & $30$ ثانیه & $15$ ثانیه \\
        گام کنترل          & $5$ میلی‌ثانیه & $0.5$ میلی‌ثانیه \\
        تعداد نمونه        & $6{,}001$ & $30{,}001$ \\
        \bottomrule
    \end{tabular}
\end{table}

\subsection{زنجیره جرم-فنر سری}
سه ارابه با جرم واحد که با فنرهای با سختی واحد به هم متصل‌اند و نیرو تنها به
ارابه نخست وارد می‌شود. کمیت مقیدشده موقعیت ارابه \emph{سوم} است، که از ورودی
با شش انتگرال‌گیر فاصله دارد. با این انتخاب، ماتریس جداساز
$L_g L_f^{5} y = k^{2}/(m_1 m_2 m_3) = 1$ به‌خوبی شرطی است و نیروهای رسم‌شده
حاصل عملکرد کنترل‌کننده‌اند و نه پیامد بدشرطی وارون‌سازی.

\subsection{پاندول دوحلقه‌ای با عملگر الاستیک}
بلوک موتور و فنر سری همان چیزی است که سامانه را به‌جای درجه نسبی $2$ به درجه
نسبی $4$ می‌رساند: گشتاور تنها از طریق فنر به حلقه می‌رسد، یعنی از میان دو
انتگرال‌گیر اضافی به ازای هر مفصل. با سختی $k=100$ نیوتن‌متر بر رادیان و اینرسی
موتور $J_m=0.1$، مُد موتور-فنر روی $\sqrt{k/J_m}=31.6$ رادیان بر ثانیه قرار
می‌گیرد؛ به‌قدر کافی بالاتر از دینامیک خودِ بازو (حدود $3$ رادیان بر ثانیه) و
به‌قدر کافی پایین‌تر از نرخ کنترل، یعنی یک رژیم واقعی عملگر الاستیک و نه یک حد
صُلب که به درجه نسبی $2$ فرو می‌کاهد.

\textbf{مأموریت پاندول کوچک نیست.} چرخش $\theta_1$ از $-\pi$ به $+\pi$ با
نگه‌داشتن $\theta_2$ در صفر خواسته شده است. هر دو نقطه ابتدایی و انتهایی، بازوی
\emph{رو به بالا} هستند و به اندازه یک دور کامل تفاوت دارند. تنها مسیر عبور از
$\theta_1=0$ می‌گذرد، جایی که بازوی باز، مجری نهایی را به $p_2^{y}=-2$ متر
می‌برد. بنابراین قید ایمنی یک کران ملایم در همسایگی مسیر نیست، بلکه در بخشی از
مأموریت با $\theta_2=0$ \emph{ناسازگار} است و بازو را وادار می‌کند برخلاف
خواسته \en{CLF} \textbf{تا شود}. همین تعارض است که وجود متغیر کمکی روی سطر
\en{CLF} و نبود آن روی سطر سد را توجیه می‌کند.

% =====================================================================
\section{نتایج}
% =====================================================================
برای هر سامانه، سه کنترل‌کننده روی یک مأموریت یکسان اجرا شده‌اند: کنترل‌کننده
مبنای \en{CLF-QP} بدون هیچ سدی، و دو اجرای \en{ECBF-CLF-QP} که در آن‌ها
\emph{سطح قید تغییر می‌کند اما قطب‌ها ثابت نگه داشته می‌شوند}. ثابت نگه‌داشتن
قطب‌ها عمدی است: به این ترتیب هر تفاوتی در پاسخ به \emph{قید} نسبت داده می‌شود
و نه به تنظیم مجدد کنترل‌کننده.

\subsection{زنجیره جرم-فنر سری (درجه نسبی ۶)}
نتایج در جدول~\ref{tab:springmass} و شکل‌های~\ref{fig:sm53}
و~\ref{fig:smbar} آمده‌اند.

\begin{table}[H]
    \centering
    \caption{نتایج حلقه‌بسته سامانه جرم-فنر سری}
    \label{tab:springmass}
    \small
    \begin{tabular}{lcccc}
        \toprule
        کنترل‌کننده & $x_3^{\max}$ & $\max x_3$ & $\min h$ & وضعیت \\
        \midrule
        \en{CLF-QP}      & $3.15$ & $3.2913$ & $-0.1413$ & نقض \\
        \en{ECBF-CLF-QP} & $3.15$ & $3.0084$ & $+0.1416$ & \textbf{ایمن} \\
        \en{ECBF-CLF-QP} & $3.00$ & $2.9624$ & $+0.0376$ & \textbf{ایمن} \\
        \midrule
        & $\max|u|$ & $p_{99}|u|$ & \% فعال & زمان نشست \\
        \midrule
        \en{CLF-QP}      & $1.314$ & $0.962$ & \nd{}   & $19.2$ \\
        \en{ECBF-CLF-QP} & $1.697$ & $1.347$ & $10.5$  & $23.4$ \\
        \en{ECBF-CLF-QP} & $1.742$ & $1.397$ & $17.6$  & $26.6$ \\
        \bottomrule
    \end{tabular}
\end{table}

دو نکته از این جدول قابل استخراج است. نخست آنکه کنترل‌کننده مبنا تا $3.2913$
فراجهش می‌کند، در حالی که مقدار منتشرشده در متن مرجع $3.27$ است؛ این نزدیکی
نشان می‌دهد بازتولید روی همان سامانه انجام شده است. دوم آنکه ادعای اصلی
شکل~\ref{fig:sm53} بازتولید می‌شود: \textbf{جابه‌جایی قید به اندازه $15$
سانتی‌متر با ثابت نگه‌داشتن قطب‌ها، نیروی بیشینه و سرعت پاسخ را عملاً بدون
تغییر می‌گذارد} ($1.70$ در برابر $1.74$ نیوتن). آنچه تغییر می‌کند سهم زمانی
فعال بودن سد است ($10.5$ در برابر $17.6$ درصد). نسبت به مبنا اما سد بها دارد: نیروی بیشینه
از $1.31$ به $1.70$ و $1.74$ نیوتن و زمان نشست از $19.2$ به $23.4$ و $26.6$ ثانیه می‌رسد.

حالت سوم شایان توجه ویژه است: با $x_3^{\max}=3.00$ هدف \emph{دقیقاً روی مرز}
قرار می‌گیرد، یعنی کنترل‌کننده باید به نقطه‌ای همگرا شود که هرگز مجاز به عبور از
آن نیست. پاسخ در $2.9624$ می‌نشیند و همان‌جا می‌ماند: ایمن، و در لحظه پایان
$3.8$ سانتی‌متر کوتاه. این بهای صریح ایمنی است.

\begin{figure*}[t]
    \centering
    \includegraphics[width=\textwidth]{ch5_springmass_fig53.png}
    \caption{پاسخ حلقه‌بسته سامانه جرم-فنر سری. ستون‌ها از چپ به راست:
    \en{CLF-QP} بدون سد، \en{ECBF-CLF-QP} با $x_3^{\max}=3.15$، و
    \en{ECBF-CLF-QP} با $x_3^{\max}=3.00$. سطر بالا موقعیت هر سه ارابه، سطر
    میانی موقعیت ارابه سوم در برابر حد مجاز (خط قرمز) و هدف (خط‌چین)، و سطر
    پایین نیروی کنترلی را نشان می‌دهد. فراجهش کنترل‌کننده مبنا در حدود
    $t\approx 5$ ثانیه از خط قرمز عبور می‌کند و چند بار از آن رد می‌شود، در حالی
    که هر دو ستون \en{ECBF} به‌صورت یکنوا از پایین به حد نزدیک می‌شوند و هرگز آن
    را لمس نمی‌کنند.}
    \label{fig:sm53}
\end{figure*}

\begin{figure*}[t]
    \centering
    \includegraphics[width=\textwidth]{ch5_springmass_barrier.png}
    \caption{بالا: مقدار تابع سد $h=x_3^{\max}-x_3$ در طول زمان. ادعای فصل آن
    است که این کمیت نامنفی بماند؛ منحنی کنترل‌کننده مبنا این ادعا را یک بار در
    $t=6.565$ ثانیه نقض می‌کند. پایین: مقدار سطری که برنامه درجه دوم واقعاً
    اعمال می‌کند، روی مقیاس لگاریتمی علامت‌دار. مقدار صفر به معنای فعال بودن قید
    است؛ لمس‌های پیاپی همان لحظاتی هستند که سد کنترل را پس می‌گیرد.}
    \label{fig:smbar}
\end{figure*}

\subsection{پاندول دوحلقه‌ای (درجه نسبی ۴)}
نتایج در جدول~\ref{tab:pendulum} و شکل‌های~\ref{fig:pend54}
و~\ref{fig:pendbar} آمده‌اند.

\begin{table}[H]
    \centering
    \caption{پاندول دوحلقه‌ای: \emph{یک} تحقق اسمی (\en{rk4}، وضعیت آغازین دقیقاً روی
    تعادل). ستون‌های پایین خواص همین تحقق‌اند و نه خواص کنترل‌کننده؛ دامنه آن‌ها روی
    گروه اجراها در جدول~\ref{tab:pendens} آمده است}
    \label{tab:pendulum}
    \small
    \begin{tabular}{lcccc}
        \toprule
        کنترل‌کننده & $p_{2\min}$ & $\min p_2^{y}$ & $\min h$ & وضعیت \\
        \midrule
        \en{CLF-QP}      & $-1.0$ & $-1.9833$ & $-0.9833$ & نقض \\
        \en{ECBF-CLF-QP} & $-1.0$ & $-0.9531$ & $+0.0469$ & \textbf{ایمن} \\
        \en{ECBF-CLF-QP} & $-0.5$ & $-0.4639$ & $+0.0361$ & \textbf{ایمن} \\
        \midrule
        & $p_{99}|\tau|$ & $\max|\tau|$ & \% فعال & $|\theta_1(T)-\pi|$ \\
        \midrule
        \en{CLF-QP}      & $26.6$ & $435.1$ & \nd{}  & $0.0544$ \\
        \en{ECBF-CLF-QP} & $22.5$ & $101.0$ & $2.8$  & $0.0013$ \\
        \en{ECBF-CLF-QP} & $22.8$ & $96.2$  & $5.4$  & $0.0152$ \\
        \bottomrule
    \end{tabular}
\end{table}

\begin{table}[H]
    \centering
    \caption{گروه اجراها: دو حل‌گر (\en{rk4}، \en{ode45}) در $9$ وضعیت آغازین
    (اسمی و هشت اغتشاش تا $2$ درجه روی $\theta(0)$)، یعنی $18$ عضو برای هر ردیف.
    دامنه‌ها تنها روی اجراهایی گرفته شده‌اند که \textbf{کامل} شده‌اند}
    \label{tab:pendens}
    \footnotesize
    \setlength{\tabcolsep}{3pt}
    \begin{tabular}{lccc}
        \toprule
        & \en{CLF-QP} & \en{ECBF}، $-1.0$ & \en{ECBF}، $-0.5$ \\
        \midrule
        کامل‌شده از $18$ & $\mathbf{18}$ & $\mathbf{6}$ & $\mathbf{4}$ \\
        \midrule
        $\min h$ & $-1.000$ تا $+0.664$ & $-0.0012$ تا $+0.664$ & $-0.0017$ تا $+0.076$ \\
        $p_{99}|\tau|$ & $13.7$ تا $42.0$ & $13.7$ تا $24.3$ & $13.0$ تا $24.6$ \\
        $\max|\tau|$ & $13.8$ تا $435.1$ & $19.6$ تا $110.5$ & $16.1$ تا $96.2$ \\
        \% فعال & \nd{} & $0.0$ تا $61.5$ & $1.5$ تا $46.0$ \\
        $|\theta_1(T)-\pi|$ & $0.0001$ تا $1.848$ & $0.0007$ تا $0.704$ & $0.0042$ تا $0.045$ \\
        $\min\theta_2$ & $-10.70$ تا $-0.09$ & $-3.67$ تا $-0.03$ & $-3.22$ تا $-2.72$ \\
        \bottomrule
    \end{tabular}
\end{table}

\textbf{هیچ‌یک از ستون‌های پایین جدول~\ref{tab:pendulum} خاصیت کنترل‌کننده نیست.}
جدول~\ref{tab:pendens} همان کمیت‌ها را روی $18$ عضو اندازه می‌گیرد و دامنه‌ها چنان
پهن‌اند که مقایسه‌های تک‌تحققی را بی‌اعتبار می‌کنند. برای نمونه، در تحقق اسمی اجرای
\en{ECBF} چرخش را \emph{دقیق‌تر} از مبنا تمام می‌کند ($0.0013$ در برابر $0.0544$)؛
روی گروه، مبنا $0.0001$ تا $1.848$ و \en{ECBF} $0.0007$ تا $0.704$ می‌دهد و دو
دامنه به‌شدت هم‌پوشانی دارند، پس آن مقایسه نتیجه‌ای عمومی نیست. به همین ترتیب قله
$435$ نیوتن‌متری مبنا لبه دامنه‌ای است که از $13.8$ آغاز می‌شود، و «سد تنها $2.8$
درصد زمان فعال است» در واقع $0$ تا $61.5$ درصد است.

\textbf{ایمنی در اجراهای کامل‌شده برقرار می‌ماند، اما کامل‌شدن شکننده است.} بدترین
$\min h$ روی اجراهای کامل‌شده \en{ECBF} برابر $-0.0012$ و $-0.0017$ است، هم‌اندازه
خروج‌های نمونه‌برداری‌شده، در حالی که مبنا به $-1.000$ می‌رسد. اینکه این دو عدد درون کران
ابعادی $c\,dt\max|\dot h|$ بخش~«\ref{sec:dt}» باشند، به $\max|\dot h|$ همان اجراها بستگی
دارد و \textbf{هنوز اندازه‌گیری نشده است}؛ تا آن زمان «خروج نمونه‌برداری‌شده» برای آن‌ها یک
فرضیه است. تقابل کیفی فصل پابرجاست. آنچه پابرجا نیست، واژه
\emph{اکیداً} است: با وضعیت آغازین اندکی جابه‌جاشده، اجرای \en{ECBF} دیگر اکیداً
ایمن نیست.

\textbf{و این هزینه‌ای دارد که تحقق اسمی آن را کاملاً پنهان می‌کند.} از $36$ اجرای
\en{ECBF} (دو سطح قید $\times$ دو حل‌گر $\times$ نه وضعیت آغازین) تنها $10$ اجرا
کامل شد: $11$ اجرا به حالت نامتناهی رسید\md{}نه مورد از آن‌ها پیش از $t=0.33$ ثانیه
و هفت مورد در بازه باریک $0.226$ تا $0.256$ ثانیه\md{}و $15$ اجرا یا حل‌گر برنامه
درجه دوم را با مقدار نامتناهی مواجه کرد ($2$) یا بودجه $180$ ثانیه‌ای دیواری را
گذراند ($13$)، که خود نشانه واگرایی است. \textbf{مبنا زیر همان اغتشاش‌ها هر $18$
اجرا را کامل می‌کند.} پس این شکنندگی به بزرگی اغتشاش برنمی‌گردد، بلکه ویژگی
کنترل‌کننده \en{ECBF} روی این سامانه است.

\textbf{سازوکار محتمل، اما تأییدنشده.} هر $18$ عضو از آزمون زنجیره شرط‌های
$y_i\ge0$ عبور می‌کنند، اما حاشیه قطب قضیه فرعی~\en{5.2} در
برخی اعضا \emph{منفی} می‌شود: کمینه $-0.395$ در سطح $-1.0$ و $-2.47$ در سطح $-0.5$.
این دو آزمون یکی نیستند و منفی‌شدن حاشیه یعنی قطب‌هایی که برای وضعیت آغازین اسمی
انتخاب شده‌اند برای وضعیت جابه‌جاشده به اندازه کافی تند نیستند. با این حال اعضایی با
حاشیه \emph{مثبت} ($+5$ و $+3.17$) نیز واگرا شدند، پس حاشیه منفی توضیح کامل نیست و
این گزارش آن را به‌عنوان علت قطعی بیان نمی‌کند.

\textbf{یک هم‌زمانی که باید آزموده شود.} هفت واگرایی از یازده در بازه $0.226$ تا $0.256$
ثانیه رخ می‌دهند، درست پس از $t=0.224$ ثانیه که قله یک‌نمونه‌ای گشتاور مبنا در تحقق اسمی
در آن رخ می‌دهد (بخش قله گشتاور کمینه\nd{}نرم): جایی که $L_gV$ از نزدیکی صفر می‌گذرد و قانون
کمینه\nd{}نرم بر $\|L_gV\|$ تقسیم می‌کند. دو توضیح پیش‌بینی‌های متفاوتی دارند. (الف) تکینگی
سطر \en{CLF} با سطر سدِ فعال برخورد می‌کند؛ آنگاه در نمونه‌های پیش از هر واگرایی، $\|L_gV\|$
کمینه، $\|\mu\|$ یا $\delta$ بیشینه و سطر سد فعال است، با هر حل‌گری. (ب) حل‌گر گام‌متغیر
\en{ode45} شکست کنترلِ نگه‌داشته‌شده را دنبال می‌کند؛ آنگاه شمار گام‌های درونی آن در هر دوره
کنترل پیش از واگرایی منفجر می‌شود\md{}همان نشانه $13$ اجرایی که بودجه دیواری را
گذراندند\md{}و \en{rk4} با گام ثابت، با یک گام در هر دوره کنترل، آنجا شکست نمی‌خورد. آزمایش
جداسازی هر $36$ عضو را با سه حل‌گر (\en{rk4} با $10$ زیرگام، \en{ode45}، و \en{rk4} با یک
گام در هر دوره) و با ثبت این کمیت‌ها در هر نمونه تکرار می‌کند، و مبنا را برای هر وضعیت
آغازین یک بار اجرا می‌کند تا زمان کمینه $\|L_gV\|$ آن معلوم شود. همان آزمایش برای هر اجرای
کامل‌شده \en{ECBF} ضریب مؤثر $c=|\min h|/(dt\max|\dot h|)$ را نیز می‌دهد. \textbf{نتیجه
هنوز اندازه‌گیری نشده است.}

این نخستین جایی است که محدودیت وابستگی قطب‌ها به وضعیت آغازین، که بخش نتیجه‌گیری
اعلام می‌کند، واقعاً \emph{آزموده} می‌شود: اجراهای اصلی این فصل همگی از حالت سکون
آغاز می‌شوند و حاشیه را با فراخی برآورده می‌کنند.

در این تحقق اسمی هر دو اجرای \en{ECBF} \textbf{اکیداً} ایمن‌اند\md{}هیچ خروجی از مرز
رخ نمی‌دهد\md{}و برنامه درجه دوم در تمام $30{,}001$ نمونه شدنی است. (جدول~\ref{tab:pendens}
نشان می‌دهد این واژه «اکیداً» تنها به همین وضعیت آغازین تعلق دارد.) نکته جالب آنکه سد تنها در
چند درصد از زمان اجرا فعال است: قطب‌های $[5,\,5.5,\,6,\,10]$ نسبت به دینامیک
بازو سریع‌اند، بنابراین سطر سد بسیار پیش از رسیدن $h$ به صفر شروع به فشار
می‌کند و مجری نهایی به حد نزدیک می‌شود و \emph{برمی‌گردد}، به‌جای آنکه روی مرز
بلغزد. این دقیقاً همان کاری است که پوش نمایی برای آن ساخته شده است. بازه فعال
بودن سد از $0.49$ تا $3.48$ ثانیه در حالت $p_{2\min}=-1.0$، به $0.49$ تا $11.19$
ثانیه در حالت تنگ‌تر گسترش می‌یابد.

\textbf{ایمنی در مختصه مقید، با تا شدن بازو در مختصه‌ای دیگر خریده می‌شود.} تابع
\en{CLF} در سراسر مأموریت $\theta_2\to 0$ را می‌خواهد، اما در اجراهای \en{ECBF}
این زاویه در میانه مأموریت تا $-2.23$ و $-3.10$ رادیان رانده می‌شود. مقدار
$-3.10$ رادیان تقریباً تا شدن کامل بازو روی خودش است، و \emph{برخورد دو حلقه با
یکدیگر در مدل وجود ندارد}. یعنی اجرای \en{ECBF} قید ارتفاع را دقیقاً برآورده
می‌کند و هم‌زمان از درون وضعیتی می‌گذرد که روی سخت‌افزار واقعی خودبرخوردی است.

این، تز خودِ این فصل است که به خودش بازمی‌گردد: سد تنها همان چیزی را تضمین می‌کند
که در $h$ نوشته شده باشد، و \textbf{قیدی که رمزگذاری نشود اعمال هم نمی‌شود}.
ایمنی در $p_2^y$ هیچ چیز درباره $\theta_2$ نمی‌گوید، و برنامه درجه دوم آزادی
باقی‌مانده را دقیقاً همان‌جا خرج می‌کند که هیچ سطری آن را نمی‌بیند. درس عملی برای
فصل ششم روشن است: هر مختصه‌ای که نقض آن هزینه فیزیکی دارد باید سطر خود را داشته
باشد.

\textbf{نوشتن قید دوم.} این درس را می‌توان آزمود. سطر دوم \en{ECBF} برای حد مفصل
$\theta_2\ge\theta_{2,\min}$ در کنار سطر ارتفاع به برنامه درجه دوم افزوده شده است، با بهره و
بررسی پذیرفتنی بودن خودش و بدون متغیر کمکی. چون $\theta$ خروجی پاندول است، این سطر پس از
خطی‌سازی دقیقاً کانال خروجی $\pm\mu_2$ است. در $\theta_1=0$ بازو برای پاک‌کردن ارتفاع باید تا
$|\theta_2|\ge\arccos(-1-p_{2\min})$ تا شود ($1.57$ رادیان در $-1.0$ و $2.09$ در $-0.5$)،
پس حدهای $-2.5$ و $-2.2$ رادیان هر دو نواری شدنی باقی می‌گذارند. آزمایش، هر دو سطح قید را با
بی‌حد، $-2.5$ و $-2.2$ رادیان روی هر دو حل‌گر اجرا می‌کند و می‌پرسد آیا دو سد درجه نسبی
چهار \emph{هم‌زمان} برقرار می‌مانند\md{}همان پرسشی که فصل ششم روی ربات راه‌رونده با آن
روبه‌رو است. \textbf{نتیجه هنوز اندازه‌گیری نشده است.}

این پیروزی سطر سختِ سد بر سطر شلِ \en{CLF} است. در همین تحقق منفرد\md{}که بخش ملاحظات نشان می‌دهد
بازتولیدپذیر نیست\md{}ایمنی این هزینه را نیز نداشت: هر دو اجرای \en{ECBF} چرخش را \emph{دقیق‌تر} از
کنترل‌کننده مبنا به پایان می‌رسانند ($0.0013$ و $0.0152$ در برابر $0.0544$
رادیان) و هر دو این کار را با گشتاور صدک نود و نهم \emph{کمتری} انجام می‌دهند.

\begin{figure*}[t]
    \centering
    \includegraphics[width=\textwidth]{ch5_pendulum_fig54.png}
    \caption{پاسخ حلقه‌بسته پاندول دوحلقه‌ای. سطر بالا دو خطای خروجی: تابع
    \en{CLF} هر دو را به صفر می‌خواهد، اما در ستون‌های \en{ECBF} خطای دوم تا
    $-2.2$ و $-3.1$ رادیان کشیده می‌شود و سپس رها می‌گردد. سطر میانی گشتاور
    مفاصل روی محور مشترک مقاوم، به‌همراه علامت‌گذاری قله واقعی. سطر پایین مسیر
    مجری نهایی در صفحه، در برابر قید (خط قرمز) و دایره دسترسی (خط‌چین). مسیر
    کنترل‌کننده مبنا در پایین نوسان مستقیماً از حد عبور می‌کند، در حالی که
    مسیرهای \en{ECBF} به‌وضوح از آن دور شده‌اند.}
    \label{fig:pend54}
\end{figure*}

\begin{figure*}[t]
    \centering
    \includegraphics[width=\textwidth]{ch5_pendulum_barrier.png}
    \caption{ارتفاع مجری نهایی در برابر حدود مجاز. کنترل‌کننده مبنا در
    $t=1.17$ ثانیه تا $-1.98$ متر سقوط می‌کند\md{}از هر دو سطح قید عبور کرده و
    تقریباً به تمام دسترسی $-2$ متری بازو می‌رسد. اجراهای \en{ECBF} به‌ترتیب در
    $-0.953$ و $-0.464$ متوقف می‌شوند، هرکدام چند سانتی‌متر پیش از خط خود، و
    بازمی‌گردند.}
    \label{fig:pendbar}
\end{figure*}

\subsection{مقایسه کنترل‌شده درجه نسبی}
این آزمایش، انگیزه اصلی فصل را به داده تبدیل می‌کند. \emph{یک} سامانه، \emph{یک}
شرط اولیه، و تنها \emph{درجه نسبی قید} تغییر می‌کند. ثابت نگه‌داشتن سامانه همان
چیزی است که این مقایسه را به شاهد تبدیل می‌کند؛ اگر سد معکوس روی سامانه دیگری
شکست می‌خورد، می‌شد آن را به گردن سامانه انداخت. نتایج در
جدول~\ref{tab:reldeg} آمده است.

\begin{table}[H]
    \centering
    \caption{بخش~\en{5.1} در برابر بخش~\en{5.2} روی یک سامانه واحد}
    \label{tab:reldeg}
    \small
    \begin{tabular}{lccc}
        \toprule
        قید & $r_b$ & $L_g h$ & کنترل‌کننده \\
        \midrule
        $\dot{x}_1 \le 0.8$ & $1$ & $-1$ & \en{CLF-QP} \\
        $\dot{x}_1 \le 0.8$ & $1$ & $-1$ & \en{CBF-CLF-QP} \\
        $x_3 \le 3.00$      & $6$ & $0$  & \en{CBF-CLF-QP} \\
        \midrule
        & قله & $\min h$ & \% ناشدنی \\
        \midrule
        $\dot{x}_1 \le 0.8$ & $0.7852$ & $+0.0148$ & $0.00$ \\
        $\dot{x}_1 \le 0.8$ & $0.4274$ & $+0.3726$ & $0.00$ \\
        $x_3 \le 3.00$      & $3.7123$ & $-0.7123$ & $\mathbf{13.50}$ \\
        \bottomrule
    \end{tabular}
\end{table}

قید با درجه نسبی یک در متن اصلی وجود ندارد و صرفاً برای این افزوده شده که بدون
آن هیچ جایی برای نشان دادن اینکه بخش~\en{5.1} \emph{کار می‌کند} باقی نمی‌ماند\md{}تنها
جاهایی که شکست می‌خورد. توجه شود که مبنای بدون سد نیز این قید را رعایت می‌کند
($\min h=+0.0148$)؛ پس سطر نخست تنها نشان می‌دهد سد درجه نسبی یک عمل می‌کند و حاشیه را
بزرگ می‌کند، نه اینکه برای این قید لازم است. \textbf{این ضعف شاهد است و باید
پذیرفته شود}: آنچه این مقایسه ثابت می‌کند تنها این است که سد معکوس در $r_b=6$
واگن می‌شود، نه اینکه در جایی که واقعاً لازم است کار می‌کند. شاهد قوی‌تر قیدی با
درجه نسبی یک است که مبنای بدون سد آن را \emph{نقض} کند: حد $\dot x_1\le0.70$، زیر قله
$0.7852$ مبنا. مقایسه با این حد نوشته شده و \textbf{هنوز اجرا نشده است}؛ یک بررسی مجموعه
آزمون همین را می‌خواهد\md{}مبنا باید حد را نقض کند و \en{ECBF} با $r_b=1$ باید آن را نگه
دارد. در اجرای با درجه نسبی شش، سطر سد تا $t=4.645$ ثانیه
بدیهی است (در حالی که $x_3=1.938$ و هنوز $h=+1.062$ با مرز فاصله دارد) و پس از
آن برای $810$ نمونه ارضاناپذیر می‌شود. هر دو حالت شکست، به همین ترتیب، مشاهده
می‌شوند.

\subsection{ملاحظات اندازه‌گیری‌شده}

\subsubsection{مسیر بازتولیدپذیر نیست، خاصیت ایمنی هست}
پس از فعال شدن سطر سد، حلقه بسته به یک \textbf{مد لغزشی زمان-گسسته} تبدیل
می‌شود: حالت میان دو نمونه از مرز خارج می‌شود و در نمونه بعد به عقب رانده
می‌گردد، و اینکه هر نمونه در کدام سمت یک تعویض مجموعه فعال بیفتد بسیار
پایین‌تر از خطای انتگرال‌گیری تعیین می‌شود. اندازه‌گیری مستقیم در
جدول~\ref{tab:integ} آمده است.

\begin{table}[H]
    \centering
    \caption{وابستگی به حل‌گر عددی، با شرط اولیه یکسان: جرم-فنر با $x_3^{\max}=3.00$ در $6$ ثانیه
    و پاندول با $p_{2\min}=-1.0$ در $3$ ثانیه. این افق‌ها کوتاه‌تر از اجراهای جدول‌های نتایج‌اند، پس
    $h_{\min}$ آن‌ها با آن جدول‌ها مقایسه‌پذیر نیست.}
    \label{tab:integ}
    \small
    \begin{tabular}{lccc}
        \toprule
        سامانه & $h_{\min}$ (\en{rk4}) & $h_{\min}$ (\en{ode45}) & واگرایی مسیر \\
        \midrule
        جرم-فنر (خطی)  & $+1.26$ & $+1.26$ & $3.2\times10^{-14}$ \\
        پاندول (غیرخطی) & $+0.104$ & $+0.0031$ & $5.6\times10^{-1}$ \\
        \bottomrule
    \end{tabular}
\end{table}

بنابراین \textbf{هیچ ادعایی در این گزارش بر یک مسیر منفرد پاندول استوار نیست}
و شکل‌ها باید به‌عنوان یک تحقق خاص خوانده شوند. آنچه هر دو حل‌گر بر آن توافق
دارند\md{}و آنچه مجموعه آزمون برای هر دو تصدیق می‌کند\md{}این است که مجموعه امن هرگز
ترک نمی‌شود.

\subsubsection{اعمال نمونه‌برداری‌شده از مرتبه \en{$O(dt)$} است: اندازه‌گیری روی سامانه خطی}
\label{sec:dt}
تغییرناپذیری پیشرو یک ادعای \emph{زمان-پیوسته} است، اما سطر مربوطه تنها در
لحظه‌های نمونه‌برداری اعمال می‌شود و سامانه میان آن‌ها حرکت می‌کند. اندازه‌گیری این
مقیاس‌بندی باید روی سامانه‌ای انجام شود که مسیرش بازتولیدپذیر باشد، وگرنه تغییرِ
اندازه‌گیری‌شده می‌تواند \emph{مسیر دیگری} باشد و نه اثر نمونه‌برداری. جدول~\ref{tab:integ}
نشان می‌دهد پاندول چنین سامانه‌ای نیست (واگرایی $5.6\times10^{-1}$)، در حالی که
جرم-فنر خطی هست ($3.2\times10^{-14}$). پس این مطالعه روی جرم-فنر انجام شده است
(جدول~\ref{tab:dt}).

\begin{table}[H]
    \centering
    \caption{مقیاس‌بندی با گام کنترل، روی جرم-فنر \emph{خطی} با $x_3^{\max}=3.00$ و
    قطب‌های ثابت. هر سطر اکیداً ایمن است؛ آنچه با $dt$ مقیاس می‌گیرد \emph{کسری حاشیه}
    $h^\star-h_{\min}$ است، که در آن $h^\star$ برون‌یابی ریچاردسون حد $dt\to0$ است}
    \label{tab:dt}
    \small
    \begin{tabular}{lccc}
        \toprule
        گام کنترل $dt$ & $h_{\min}$ & کسری حاشیه & نسبت \\
        \midrule
        $4$ میلی‌ثانیه     & $+0.0376392$ & $1.468\times10^{-4}$ & \nd{} \\
        $2$ میلی‌ثانیه     & $+0.0377122$ & $7.374\times10^{-5}$ & $1.98$ \\
        $1$ میلی‌ثانیه     & $+0.0377490$ & $3.693\times10^{-5}$ & $2.00$ \\
        $0.5$ میلی‌ثانیه   & $+0.0377675$ & $1.849\times10^{-5}$ & $1.99$ \\
        $0.25$ میلی‌ثانیه  & $+0.0377767$ & $9.246\times10^{-6}$ & $2.00$ \\
        $0.125$ میلی‌ثانیه & $+0.0377813$ & $4.623\times10^{-6}$ & $2.00$ \\
        \bottomrule
    \end{tabular}
\end{table}

شیب لگاریتمی کسری حاشیه در برابر $dt$ روی این شش نقطه، که بازه‌ای $32$ برابری را
می‌پوشانند، برابر $\mathbf{0.998}$ است و هر نصف‌کردن $dt$ کسری را دقیقاً نصف می‌کند
(نسبت‌های $1.98$ تا $2.00$). \textbf{این ادعای $O(dt)$ را اندازه‌گیری می‌کند و نه
مفروض می‌گیرد}، و دو نکته را روشن می‌سازد. نخست آنکه روی این سامانه در هیچ گام
کنترلی خروجی از مرز رخ نمی‌دهد؛ آنچه نمونه‌برداری هزینه می‌کند \emph{حاشیه} است و نه
\emph{ایمنی}. دوم آنکه روی پاندول \textbf{هیچ مقیاس‌بندی تمیزی مشاهده نشد}: سه گام کنترل
نتیجه‌ای با تغییر علامت دادند\md{}از $-1.2\times10^{-3}$ به $+4.7\times10^{-2}$ در $0.5$
میلی‌ثانیه\md{}که مسیری دیگر است و نه خروجی کوچک‌تر، و این گزارش ادعای $O(dt)$ را از روی
آن بیان نمی‌کند.

\textbf{کران ابعادی خروج.} برچسب «خروج نمونه‌برداری‌شده» تنها وقتی به یک $\min h$ منفی داده
می‌شود که $|\min h|\le c\,dt\max|\dot h|$، با $c=10$ و $\max|\dot h|$ اندازه‌گیری‌شده از
تاریخچه $h$ همان اجرا؛ این کران از نظر ابعادی یک طول است. روی جرم-فنر
$\max|\dot h|=0.665$ است، پس کران در $dt=1$ میلی‌ثانیه برابر $6.65\times10^{-3}$ می‌شود.
روی گروه اجراهای پاندول این کران هنوز اعمال نشده است (بخش پاندول).

\subsubsection{قله گشتاور کمینه-نرم}
یک \en{CLF-QP} کمینه-نرم مقدار $\|\mu\| = \psi/\|L_g V\|$ را اعمال می‌کند که
هرجا $L_g V$ از نزدیکی صفر عبور کند بی‌کران می‌شود. روی کنترل‌کننده مبنای
پاندول این اتفاق یک بار در $t=0.224$ ثانیه رخ می‌دهد، جایی که گشتاور برای
\textbf{یک نمونه از $30{,}001$ نمونه} به $435$ نیوتن‌متر می‌رسد، در برابر صدک
نود و نهم برابر با $26.6$. این ویژگی صورت‌بندی کمینه-نرم است و نه ویژگی سد؛
اجراهای \en{ECBF} چنین نمونه‌ای ندارند. به همین دلیل $\max|u|$ همواره در کنار
$p_{99}|u|$ گزارش می‌شود.

\subsubsection{نکته‌ای درباره حل‌گر برنامه درجه دوم}
حل‌گر در این فصل روش \emph{مجموعه فعال} را اجرا می‌کند و نه روش \emph{نقطه
درونی} که فصل‌های سوم و چهارم از آن استفاده می‌کنند، و این تفاوت حاشیه‌ای نیست.
هنگامی که سطر سد به‌شدت فشار می‌آورد، تابع هدف \en{CLF} و قید \en{ECBF} با
اختلاف زیادی مخالف یکدیگرند، و روش نقطه درونی روی چنین مسئله‌هایی می‌تواند با کد
خروج $-3$ («بی‌کران») بازگردد، آن هم روی مسئله‌ای که هسیَن آن مثبت معین است و
بنابراین \emph{نمی‌تواند} بی‌کران باشد. به همین دلیل این فصل از روش مجموعه فعال
استفاده می‌کند. این انتخاب بر پایه همین استدلال ساختاری است و نه بر پایه یک شمارش: شمار
چنین شکست‌هایی در اجراهای این گزارش اندازه‌گیری نشده است.

\subsubsection{آغاز از $x_0$ متحرک: محدودیت وابستگی به شرط اولیه، اندازه‌گیری‌شده}
\label{sec:moving}
همه اجراهای جدول‌های نتایج از سکون آغاز می‌شوند و از این رو شرط پذیرفتنی بودن قطب‌ها
را با حاشیه کامل برآورده می‌کنند (جدول~\ref{tab:adm}). آزمایش $x_0$ متحرک همان
چهار پیکربندی را با سرعتی اولیه به سمت مرز آغاز می‌کند: در جرم\nd{}فنر هر سه ارابه
با هم با سرعت $s$، و در آونگ حلقه نخست و موتورش با هم با سرعت زاویه‌ای $s$ از همان
وضعیت رو به بالا. برای هر پیکربندی ابتدا سرعتی که زنجیره در آن می‌شکند با دوبخشی
یافته می‌شود و سپس حلقه بسته در مضرب‌هایی از آن اجرا می‌گردد
(جدول~\ref{tab:moving}).

\begin{table}[H]
    \centering
    \caption{آغاز از $x_0$ متحرک. $s_{\mathrm{adm}}$ سرعتی است که پله‌های دقیقاً
    بررسی‌پذیر زنجیره در آن می‌شکنند و $s_{\mathrm{top}}$ سرعتی که $y_{r_b}$ بدون کنترل
    منفی می‌شود (متر بر ثانیه برای جرم\nd{}فنر، رادیان بر ثانیه برای آونگ). مضرب‌های
    آزموده‌شده از $s_{\mathrm{adm}}$: $0.5$، $0.9$، $1.1$، $1.5$، $2$ و $3$. ستون آخر
    اجراهایی است که همه قطب‌ها در $1.05$ برابر کوچک‌ترین ضریبِ پذیرفتنی‌کننده $x_0$
    ضرب شده‌اند}
    \label{tab:moving}
    \footnotesize
    \begin{tabular}{lccccc}
        \toprule
        اجرا & $s_{\mathrm{adm}}$ & $s_{\mathrm{top}}$ & ایمن تا & نخستین نقض & قطب‌های سریع‌تر \\
        \midrule
        جرم\nd{}فنر، $3.15$ & $0.894$ & $0.773$ & $\times1.5$ & $\times2$: $-0.162$ & ایمن \\
        جرم\nd{}فنر، $3.00$ & $0.852$ & $0.736$ & $\times1.5$ & $\times2$: $-0.159$ & ایمن \\
        آونگ، $-1.0$ & $3.73$ & $3.06$ & $\times0.5$ & $\times0.9$: واگرا & واگرا \\
        آونگ، $-0.5$ & $3.41$ & $2.78$ & \nd{} & $\times0.5$: واگرا & واگرا \\
        \bottomrule
    \end{tabular}
\end{table}

\textbf{روی جرم\nd{}فنر، کران پذیرفتنی بودن محافظه‌کارانه است و قطب‌های سریع‌تر آن را
جبران می‌کنند.} دوبخشی همان کران شکل بسته را می‌دهد (اختلاف $10^{-16}$) و پله پنجم
است که می‌بندد. اما $h$ تا $1.5$ برابر $s_{\mathrm{adm}}$ ($1.34$ و $1.28$ متر بر
ثانیه) نامنفی می‌ماند و تنها در $2$ برابر ($1.79$ و $1.70$) به $-0.162$ و $-0.159$ و در
$3$ برابر به $-1.48$ و $-1.41$ می‌رسد. پس شرط کافی است و نه لازم، و روی این سامانه با
ضریبی میان $1.5$ و $2$ در سرعت محافظه‌کار است. قطب‌های سریع‌تر $h$ را در هر سه سرعتی
که آزموده شدند ($1.1$، $1.5$ و $3$ برابر) نامنفی نگه می‌دارند، اما با هزینه: در $3$ برابر $s_{\mathrm{adm}}$ قله نیرو از
$84$ و $80$ نیوتن به $413{,}010$ و $71{,}066$ نیوتن می‌رسد، زیرا در درجه نسبی $6$
بهره‌های سد تا توان ششم ضریب قطب‌ها رشد می‌کنند.

\textbf{روی آونگ، پیش از آنکه این شرط اهمیت یابد، خودِ کنترل‌کننده از کار می‌افتد.}
تنها کندترین آغاز متحرک، $1.87$ رادیان بر ثانیه در $p_{2\min}=-1.0$، کامل می‌شود
($h_{\min}=+0.52$). بقیه\md{}از $0.9\,s_{\mathrm{adm}}$ در $-1.0$ و از
$0.5\,s_{\mathrm{adm}}$ ($1.70$ رادیان بر ثانیه) در $-0.5$، پذیرفتنی یا نه، و همه
اجراهای با قطب‌های سریع‌تر\md{}ظرف $0.13$ تا $1.17$ ثانیه به حالت نامتناهی می‌رسند،
با برنامه درجه دوم ناشدنی و گشتاورهایی بالای $10^{11}$ نیوتن‌متر؛ $15$ اجرا از این
$17$ پیش از واگرایی از مرز می‌گذرند. علت، دست‌کم بالای $s_{\mathrm{top}}$، ساختاری
است: ضریب کنترل سطر سد، $L_gL_f^{3}h=(k/J_m)\,\nabla p_2^{y\,\top}D^{-1}$، هرجا
$\nabla p_2^y=0$ باشد صفر است، یعنی هرگاه بازو قائم باشد\md{}و وضعیت آغاز مأموریت دقیقاً
چنین است. پس درجه نسبی قید در $x_0$ برابر $4$ نیست و فرض قضیه~$5.1$ همان‌جا برقرار
نیست. از سکون این بی‌خطر است، چون سطر در $t=0$ شل است ($y_4=4950$ و $4120$)؛ اما
بالای $s_{\mathrm{top}}$ مقدار $y_4(x_0)$ منفی است و چون کنترل در $x_0$ در آن ظاهر
نمی‌شود، برنامه درجه دوم از نخستین نمونه ناشدنی است. اجرای $1.70$ رادیان بر ثانیه در
$-0.5$، که پذیرفتنی است و سطرش در $t=0$ شل، دیرتر ($t=1.17$ ثانیه) از کار می‌افتد.
بازوی کاملاً تاشده با حلقه نخست رو به پایین ($\theta_1=0$، $\theta_2=\pm\pi$) صفر دیگری از
همین ضریب است و اجراهای اسمی \en{ECBF} خود $\theta_2$ را تا $-3.10$ رادیان تا می‌کنند؛
این نامزد محتملی است، اما این مطالعه علت را مکان‌یابی نمی‌کند.

\subsection{اعتبارسنجی نرم‌افزاری}
مجموعه آزمون در چهار لایه سازمان یافته است و ترتیب آن عمدی است: هر لایه لایه
زیرین خود را مفروض می‌گیرد. نتایج در جدول~\ref{tab:tests} آمده است. آخرین اجرای
اندازه‌گیری‌شده $61$ بررسی داشت که همه در $51.2$ ثانیه قبول شدند. افزون بر این‌ها دو بررسی
تازه برای سد حد مفصل هنوز اجرا نشده‌اند (سطر آن دقیقاً $\pm\mu_2$ است، و در برنامه درجه دوم
کنار سطر ارتفاع می‌نشیند و در وضعیت آغاز پذیرفتنی است)، و معیار یک بررسی موجود سخت‌تر شده
است: با حد $0.70$، سد درجه نسبی یک باید قیدی را نگه دارد که مبنا \emph{نقض} می‌کند.

\begin{table}[H]
    \centering
    \caption{مجموعه آزمون فصل پنجم: آخرین اجرای اندازه‌گیری‌شده ($61$ بررسی، $51.2$ ثانیه)
    و بررسی‌هایی که هنوز اجرا نشده‌اند}
    \label{tab:tests}
    \small
    \begin{tabular}{lccc}
        \toprule
        پیمانه & بررسی & نتیجه & اجرانشده \\
        \midrule
        مدل دو سامانه و درجه‌های نسبی & $15$ & قبول & \nd{} \\
        پشته مشتق‌های لی و تابع سد    & $15$ & قبول & \nd{} \\
        نظریه بخش~\en{5.2} و تضمین ایمنی   & $17$ & قبول & $2$ \\
        برنامه‌های درجه دوم           & $14$ & قبول & $1$ تغییریافته \\
        \midrule
        \textbf{مجموع} & $\mathbf{61}$ & \textbf{قبول} & $\mathbf{2}+1$ \\
        \bottomrule
    \end{tabular}
\end{table}

دو بررسی به‌طور ویژه شایان ذکرند، زیرا در برابر \emph{نادرستی خاموش} محافظت
می‌کنند و نه نادرستی آشکار. نخست، تصدیق $\mathrm{eig}(F_b-G_bK_b)=-p$ که خطای
معکوس‌شدن ترتیب بهره را می‌گیرد\md{}خطایی که در هر دو حالت ماتریسی هورویتس تولید
می‌کند و بنابراین هیچ جای دیگری اعتراضی برنمی‌انگیزد. دوم، مشتق‌های مرتبه چهارم
سد که به \emph{دو روش مستقل} اعتبارسنجی می‌شوند: در برابر تفاضل‌های متناهی
(نه‌چندان دقیق، اما حساس به یک مشتق نادرست) و در برابر یک زنجیره قاعده
فاآ دی برونو نوشته‌شده به‌دست (دقیق، تا $1.4\times10^{-15}$). دلیل این
سخت‌گیری آن است که تفاضل‌های متناهی تودرتو در هر مرتبه حدود سه رقم دقت از دست
می‌دهند و در مرتبه چهارم خطا به $O(10^{-3})$ می‌رسد\md{}روی کمیتی که کنترل‌کننده
دقیقاً روی آن می‌نشیند.

% =====================================================================
\section{ادعاها و وضعیت آن‌ها}
% =====================================================================
جدول~\ref{tab:claims5} هر ادعای این گزارش را با شاهد آن و وضعیتش کنار هم می‌گذارد.
«اندازه‌گیری‌شده» یعنی عدد از اجرایی است که در جدول بازتولید (پیوست~\ref{app:repro5}) نام
برده شده؛ «در انتظار اجرا» یعنی آزمایش نوشته شده اما هنوز اجرا نشده، و گزارش تا آن زمان
چیزی درباره نتیجه‌اش ادعا نمی‌کند.

\begin{table*}[!t]
    \centering
    \caption{ادعاهای فصل پنجم}
    \label{tab:claims5}
    \small
    \begin{tabular}{p{0.33\textwidth}p{0.43\textwidth}p{0.14\textwidth}}
        \toprule
        ادعا & شاهد & وضعیت \\
        \midrule
        سد معکوس تنها در درجه نسبی یک قید مؤثری بر ورودی می‌سازد &
            $L_gh\equiv0$ در $r_b=6$؛ سطر بدیهی تا $t=4.645$ و سپس ارضاناپذیر در $13.5$ درصد نمونه‌ها &
            اندازه‌گیری‌شده \\
        سد درجه نسبی یک آنجا که لازم است کار می‌کند &
            حد $\dot x_1\le0.70$، زیر قله $0.7852$ مبنا &
            در انتظار اجرا \\
        \en{ECBF} در درجه نسبی $6$ و $4$ ایمنی را نگه می‌دارد (تحقق‌های اسمی) &
            $h_{\min}=+0.14$، $+0.038$، $+0.047$ و $+0.036$؛ برنامه درجه دوم در $72{,}004$ نمونه شدنی &
            اندازه‌گیری‌شده \\
        با اغتشاش $2$ درجه، اجراهای کامل‌شده ایمن‌اند اما نه «اکیداً» &
            بدترین $\min h=-0.0017$؛ تنها $10$ از $36$ اجرا کامل &
            اندازه‌گیری‌شده \\
        $-0.0017$ درون کران ابعادی خروج نمونه‌برداری‌شده است &
            ضریب مؤثر $c$ روی اجراهای کامل‌شده &
            در انتظار اجرا \\
        واگرایی‌ها از تکینگی \en{CLF} در کنار سد فعال‌اند یا از حل‌گر گام‌متغیر &
            سه حل‌گر و ثبت $\|L_gV\|$، $\|\mu\|$، $\delta$ و گام‌های درونی &
            در انتظار اجرا \\
        نمونه‌برداری حاشیه را به اندازه $O(dt)$ هزینه می‌کند و نه ایمنی را &
            جرم\nd{}فنر: شیب $0.998$ روی بازه $32$ برابری $dt$؛ روی پاندول مقیاس‌بندی تمیزی دیده نشد &
            اندازه‌گیری‌شده (تنها خطی) \\
        شرط پذیرفتنی بودن در $x_0$ کافی است و نه لازم &
            جرم\nd{}فنر: ایمن تا $1.5\,s_{\mathrm{adm}}$، نقض در $2\,s_{\mathrm{adm}}$ &
            اندازه‌گیری‌شده \\
        از $x_0$ متحرک، سد پاندول در بازوی قائم درجه نسبی خود را از دست می‌دهد &
            $\nabla p_2^y=0$ در وضعیت آغاز؛ $17$ از $18$ اجرا واگرا &
            اندازه‌گیری‌شده؛ یک اجرا مکان‌یابی‌نشده \\
        قیدی که نوشته نشود اعمال نمی‌شود &
            $\theta_2$ تا $-3.10$ رادیان برای برآوردن قید ارتفاع &
            اندازه‌گیری‌شده \\
        دو سد درجه نسبی چهار (ارتفاع و حد مفصل) هم‌زمان برقرار می‌مانند &
            دو سطح قید، سه حالت حد، دو حل‌گر &
            در انتظار اجرا \\
        حاشیه قضیه فرعی~\en{5.2} در حالت پس از برخورد گام مرجع &
            سدهای جای پای فصل ششم در نقطه ثابت گام مرجع &
            در انتظار اجرا \\
        \bottomrule
    \end{tabular}
\end{table*}

% =====================================================================
\section{نتیجه‌گیری}
% =====================================================================
رشته‌ای که این فصل را به فصل‌های سوم و چهارم می‌بندد این است: \emph{تضمینی که در زمان
پیوسته اثبات شده، در پیاده‌سازیِ نمونه‌برداری‌شده با گشتاور محدود و تماس محدود دوام
نمی‌آورد\md{}و هر تضمین تنها آنچه را می‌پوشاند که در آن نوشته شده است.} در این فصل هر دو
نیمه آن دیده می‌شود: تغییرناپذیری پیشرو در زمان پیوسته اثبات می‌شود و اعمال نمونه‌برداری‌شده
از حاشیه آن می‌کاهد، و سدی که ارتفاع را بی‌نقص نگه می‌دارد، مفصلی را که در $h$ نوشته نشده
تا حد خودبرخوردی رها می‌کند.

یک تابع سد کنترلی نمایی سطری است که نسبت به ورودی مجازی اَفین است، در هر درجه
نسبی کار می‌کند، با جایابی قطب طراحی می‌شود، و در همان برنامه درجه دومی می‌نشیند
که سطر پایداری در آن قرار دارد. شش شبیه‌سازی حلقه‌بسته اسمی روی دو سامانه با
دینامیک، تعداد ورودی و درجه نسبی متفاوت نشان دادند که کنترل‌کننده مبنا در هر دو
حالت مجموعه امن را ترک می‌کند، در حالی که تمام اجراهای \en{ECBF} ایمن باقی
می‌مانند و برنامه درجه دوم در هیچ نمونه‌ای ناشدنی نمی‌شود؛ با اغتشاش کوچکی در وضعیت آغازین
پاندول اما تنها $10$ از $36$ اجرای \en{ECBF} کامل می‌شود. همچنین نشان داده شد
که ایمنی لزوماً به بهای کیفیت ردیابی تمام نمی‌شود: در تحقق گزارش‌شده پاندول، اجراهای ایمن
مأموریت را دقیق‌تر و با گشتاور صدک نود و نهم کمتری به پایان رساندند. چون مسیر پاندول
بازتولیدپذیر نیست، این یک مشاهده است و نه نتیجه‌ای عمومی؛ روی جرم-فنر، سد نیروی بیشینه و
زمان نشست را افزایش داد.

آنچه این فصل به فصل ششم تحویل می‌دهد، همین ساختار است؛ فصل ششم آن را بر راه‌رفتن
با طول و عرض گام متغیر و بر سنگ‌های پله متحرک اعمال می‌کند، جایی که «اصلاح سد
برای قیدهای مبتنی بر موقعیت» چیزی جز حالت درجه نسبی $2$ همین ساختار نیست
\cite{nguyen2015footstep}. این یکسانی یک تشبیه نیست و در هر دو سو به‌صورت عددی
تصدیق شده است.

دو محدودیت باید همراه این ساختار حمل شوند. نخست آنکه \textbf{انتخاب قطب‌ها به
شرط اولیه وابسته است}؛ مجموعه‌ای از قطب‌ها که برای یک شرط اولیه معتبر است،
خودبه‌خود برای شرط دیگری معتبر نیست. روی یک ربات راه‌رونده، که در آن $x_0$ هرچه
باشد خروجی نگاشت ضربه گام پیشین است، این یک قید زنده است و نه یک تشریفات. از
$x_0$ متحرک (بخش~\ref{sec:moving})، این کران روی جرم\nd{}فنر با ضریب $1.5$ تا $2$
محافظه‌کار است و قطب‌های سریع‌تر آن را به بهای نیرو جبران می‌کنند؛ روی آونگ، آغاز
متحرک پیش از آن به وضعیتی می‌خورد که سد در آن درجه نسبی خود را از دست می‌دهد، و این
درسی است برای هر سدی روی ارتفاع یک نقطه در فصل ششم. پیش‌نمایش همین محدودیت روی راه‌رفتن
نوشته شده است: شرط قضیه فرعی~\en{5.2} در حالت پس از برخورد نقطه ثابت گام مرجع، با سدهای
جای پای فصل ششم و قطب‌هایی که آن فصل به کار می‌برد\md{}کمینه قطب پذیرفتنی
$\gamma_{b,\min}=-\dot g/g$ در برابر قطب در حال استفاده $\gamma_b$، یعنی حاشیه‌ای که هر
برخورد مختل‌شده از آن آغاز می‌کند. \textbf{نتیجه آن هنوز اندازه‌گیری نشده است.} دوم
آنکه \textbf{تمام این ساختار مبتنی بر مدل است}: هر جمله در سطر سد از مدل
می‌آید. این دقیقاً همان شکافی است که فصل چهارم تمام حجم خود را صرف آن کرد و
هیچ‌چیز در فصل حاضر به آن نمی‌پردازد \cite{xu2015robustness}. ترکیب یک سد نمایی
مقاوم یا تطبیقی، گام طبیعی بعدی است که در اینجا آزموده نشده است.

سوم، هر دو سامانه اعتبارسنجی \emph{فاقد دینامیک صفر} هستند؛ رابطه
$r\cdot n_y = n_x$ در هر دو برقرار است و هیچ چیزی نمی‌تواند در یک مُد داخلی
پنهان شود. این یک راحتی مربوط به همین سامانه‌هاست که ربات راه‌رونده فصل ششم از
آن برخوردار نخواهد بود \cite{westervelt2007feedback}.

سرانجام، روشن‌ترین درس این فصل از خودِ نتیجه موفق آن بیرون می‌آید. اجراهای
\en{ECBF} قید ارتفاع مجری نهایی را بی‌نقص اعمال می‌کنند و برای این کار
$\theta_2$ را تا $-3.10$ رادیان، یعنی تقریباً تا شدن کامل بازو روی خودش، پیش
می‌برند\md{}وضعیتی که روی سخت‌افزار واقعی خودبرخوردی است و \emph{در مدل وجود
ندارد}. سد دقیقاً همان چیزی را تضمین کرد که در $h$ نوشته شده بود، و آزادی
باقی‌مانده را همان‌جا خرج کرد که هیچ سطری آن را نمی‌دید. \textbf{قیدی که
رمزگذاری نشود اعمال هم نمی‌شود}؛ ایمنی یک خاصیت مدل است و نه خاصیت ربات، و
گستره آن دقیقاً به گستره $h$ محدود است. برای فصل ششم این یعنی هر مختصه‌ای که
نقض آن هزینه فیزیکی دارد\md{}خودبرخوردی، حدود مفصلی، مخروط اصطکاک\md{}باید سطر
خود را داشته باشد، وگرنه برنامه درجه دوم آن را به‌عنوان آزادی رایگان مصرف
خواهد کرد.

% ---------------------- مراجع ----------------------
\begin{LTR}
\printbibliography[title={مراجع}]
\end{LTR}


% =====================================================================
% پیوست: استخراج روابط
% =====================================================================
\onecolumn
\setcounter{section}{0}
\renewcommand{\thesection}{پ}
\renewcommand{\thesubsection}{پ-\arabic{subsection}}
\section{استخراج روابط}
\label{app:deriv}
این پیوست هر رابطه متن را گام‌به‌گام به دست می‌آورد و ترتیب آن همان ترتیب متن است.
نمادگذاری: $r_b$ درجه نسبی قید، $\eta_b=(h,\dot h,\ldots,h^{(r_b-1)})$، و
$F_b$، $G_b$ زوج کانونی زنجیره $r_b$ انتگرال‌گیر (یک روی قطر بالایی، $G_b=e_{r_b}$).

% ---------------------------------------------------------------------
\subsection{سد معکوس (روابط~\eqref{eq:recip} و \eqref{eq:degenerate})}
با $B=1/h$ داریم $\dot B=-\dot h/h^{2}$. شرط $\dot B\le\gamma/B=\gamma h$ را بنویسید:
\begin{equation*}
    -\frac{\dot h}{h^{2}}\le\gamma h
    \quad\Longrightarrow\quad
    \dot h\ge-\gamma h^{3}
\end{equation*}
(ضرب در $h^2>0$ جهت نامساوی را حفظ می‌کند، و این دقیقاً جایی است که فرض $h>0$ لازم است).
با $\dot h=L_fh+L_gh\,u$ سطر $L_gh\,u\ge-\gamma h^3-L_fh$ به دست می‌آید. اگر درجه نسبی
$r_b>1$ باشد، بنا به تعریف $L_gh\equiv0$ و سمت چپ صفر می‌شود؛ رابطه~\eqref{eq:degenerate}.

% ---------------------------------------------------------------------
\subsection{سطر نمایی اَفین در $\mu$ (روابط~\eqref{eq:viol} و \eqref{eq:chain})}
درجه نسبی $r_b$ یعنی $L_gL_f^{i}h\equiv0$ برای $i<r_b-1$. پس مشتق‌های پیاپی تا مرتبه
$r_b-1$ ورودی ندارند، $h^{(i)}=L_f^{i}h$، و در مشتق مرتبه $r_b$ ورودی نخستین بار ظاهر می‌شود:
\begin{equation*}
    h^{(r_b)}=\frac{\partial L_f^{r_b-1}h}{\partial x}\,(f+gu)=L_f^{r_b}h+L_gL_f^{r_b-1}h\;u .
\end{equation*}
با جای‌گذاری $u=u_{\mathrm{ff}}+A^{-1}\mu$:
\begin{equation*}
    h^{(r_b)}=\underbrace{L_f^{r_b}h+L_gL_f^{r_b-1}h\,u_{\mathrm{ff}}}_{b_0(x)}
    +\underbrace{L_gL_f^{r_b-1}h\,A^{-1}}_{L_b(x)}\,\mu=:\mu_b ,
\end{equation*}
که در $\mu$ اَفین است. چون مؤلفه $i$ام $\eta_b$ مشتق مؤلفه $(i-1)$ام است و آخرین مشتق $\mu_b$،
$\dot\eta_b=F_b\eta_b+G_b\mu_b$.

% ---------------------------------------------------------------------
\subsection{جایابی قطب و ترتیب وارون بهره‌ها}
$A_b=F_b-G_bK_b$ ماتریس همراهی است که سطر آخر آن $-K_b=-[k_1,\ldots,k_{r_b}]$ است. چندجمله‌ای
مشخصه آن
\begin{equation*}
    \det(sI-A_b)=s^{r_b}+k_{r_b}s^{r_b-1}+\cdots+k_2s+k_1 .
\end{equation*}
برای قطب‌های $-p_1,\ldots,-p_{r_b}$ آن را با
$\prod_i(s+p_i)=s^{r_b}+a_1s^{r_b-1}+\cdots+a_{r_b}$ برابر کنید: $k_j=a_{r_b+1-j}$، یعنی
$K_b=[a_{r_b},\ldots,a_1]$. ضرایب به ترتیب \emph{وارون} می‌آیند؛ ترتیب مستقیم نیز چندجمله‌ای
هورویتس می‌دهد، اما با قطب‌های دیگر.

% ---------------------------------------------------------------------
\subsection{از قانون کنترل به پوش پایینی (رابطه~\eqref{eq:ecbf})}
\textbf{راه اول: پاسخ ضربه.} شکاف را $\nu:=\mu_b+K_b\eta_b\ge0$ بنامید. آنگاه
\begin{equation*}
    \dot\eta_b=(F_b-G_bK_b)\eta_b+G_b\nu=A_b\eta_b+G_b\nu ,
\end{equation*}
و با $C_b=e_1^{\top}$ (که $h$ را برمی‌گزیند):
\begin{equation*}
    h(t)=C_be^{A_bt}\eta_b(0)+\int_0^t\underbrace{C_be^{A_b(t-\tau)}G_b}_{g(t-\tau)}\,\nu(\tau)\,d\tau .
\end{equation*}
$g$ پاسخ ضربه تابع تبدیل $1/\prod_i(s+p_i)$ است، یعنی کانولوشن $e^{-p_1t},\ldots,e^{-p_{r_b}t}$. اگر
همه $p_i$ حقیقی و مثبت باشند، هر عامل نامنفی است و کانولوشن نیز، پس $g\ge0$. با $\nu\ge0$ انتگرال
نامنفی است و
\begin{equation*}
    h(t)\ge C_be^{A_bt}\eta_b(0).
\end{equation*}
\textbf{چرا قطب‌ها باید حقیقی باشند.} برای یک زوج مزدوج مختلط پاسخ ضربه نوسان می‌کند و علامت
عوض می‌کند؛ آنگاه $\nu\ge0$ می‌تواند $h$ را \emph{کاهش} دهد و پوش پایینی از میان می‌رود، هرچند
$A_b$ هورویتس است.

\textbf{راه دوم: زنجیره مجموعه‌ها (استدلال قضیه).} تعریف کنید
\begin{equation*}
    y_0=h,\qquad y_i=\dot y_{i-1}+p_iy_{i-1},\quad i=1,\ldots,r_b .
\end{equation*}
با بسط، $y_{r_b}=\prod_i(\tfrac{d}{dt}+p_i)\,h=h^{(r_b)}+a_1h^{(r_b-1)}+\cdots+a_{r_b}h=\mu_b+K_b\eta_b$
(بخش پ-۳). پس سطر ایمنی دقیقاً $y_{r_b}\ge0$ است. حال از بالا به پایین:
اگر $y_i\ge0$، آنگاه $\dot y_{i-1}\ge-p_iy_{i-1}$ و طبق لم مقایسه
$y_{i-1}(t)\ge e^{-p_it}y_{i-1}(0)$. اگر $y_{i-1}(0)\ge0$ آنگاه $y_{i-1}(t)\ge0$ برای همه $t$. با
تکرار تا $i=1$ نتیجه $h(t)\ge0$ است. هر پله به عدد حقیقی مثبت $p_i$ نیاز دارد؛ با قطب مختلط،
$y_i$ حقیقی تعریف نمی‌شود.

% ---------------------------------------------------------------------
\subsection{شرط پذیرفتنی بودن قطب‌ها در $x_0$}
استدلال راه دوم \emph{همه} شرط‌های $y_{i-1}(0)\ge0$ را لازم دارد، نه فقط $h(0)\ge0$. از
$y_i(0)=\dot y_{i-1}(0)+p_iy_{i-1}(0)\ge0$ و $y_{i-1}(0)>0$:
\begin{equation*}
    p_i\ \ge\ -\frac{\dot y_{i-1}(x_0)}{y_{i-1}(x_0)},\qquad i=1,\ldots,r_b .
\end{equation*}
برای $i=1$ این یعنی $p_1\ge-\dot h(0)/h(0)$: اگر $h$ در لحظه آغاز با سرعت $\dot h(0)<0$ به سمت مرز
می‌رود، قطب باید دست‌کم به همان نسبت سریع باشد. در پله آخر، $y_{r_b}(0)$ به $h^{(r_b)}$ و در نتیجه به
کنترل وابسته است؛ پیاده‌سازی آن را در مقدار رانش ($u=0$) می‌سنجد.

% ---------------------------------------------------------------------
\subsection{برنامه درجه دوم (رابطه~\eqref{eq:qp})}
هر سطر نسبت به $\mu$ اَفین است، پس مسئله یک \en{QP} واقعی است:
\begin{itemize}\renewcommand{\labelitemi}{\lr{\textbullet}}
    \item پایداری: $\dot V+\lambda V=L_fV+L_gV\mu+\lambda V\le\delta$.
    \item ایمنی: $b_0+L_b\mu\ge-K_b\eta_b$، یعنی $-L_b\mu\le b_0+K_b\eta_b$؛ بدون متغیر کمکی.
    \item قید ورودی: $u=u_{\mathrm{ff}}+A^{-1}\mu$ در $\mu$ اَفین است، پس جعبه $|u|\le u_{\max}$
          به سطرهای $A_c\mu\le b_c$ تبدیل می‌شود.
\end{itemize}
با $\delta$ آزاد، سطر پایداری همیشه برقرارشدنی است؛ ناشدنی بودن تنها می‌تواند از سطر ایمنی بیاید:
یا در تعارض با قید ورودی، یا با $L_b=0$ و سمت راست نامطلوب.

% ---------------------------------------------------------------------
\subsection{درجه نسبی دو سامانه اعتبارسنجی}
\textbf{جرم\nd{}فنر.} معادلات:
\begin{equation*}
    m_1\ddot x_1=u-k(x_1-x_2),\quad
    m_2\ddot x_2=k(x_1-x_2)-k(x_2-x_3),\quad
    m_3\ddot x_3=k(x_2-x_3).
\end{equation*}
خروجی $x_3$ را مشتق بگیرید و تنها جمله‌ای را دنبال کنید که به $u$ نزدیک‌تر می‌شود:
$\ddot x_3\ni\tfrac{k}{m_3}x_2$؛ $x_3^{(4)}\ni\tfrac{k}{m_3}\ddot x_2\ni\tfrac{k^2}{m_2m_3}x_1$؛
$x_3^{(6)}\ni\tfrac{k^2}{m_2m_3}\ddot x_1\ni\tfrac{k^2}{m_1m_2m_3}u$. هیچ مشتقی پیش از مرتبه $6$ جمله‌ای از $u$
ندارد، پس درجه نسبی $6$ و $L_gL_f^5y=k^2/(m_1m_2m_3)=1$.

\textbf{پاندول با عملگر الاستیک.} گشتاور مفصل $u_{\mathrm{joint}}=-k(\theta-\theta_m)-\xi\dot\theta$ و
موتور $J_m\ddot\theta_m=k(\theta-\theta_m)+\tau$. در $\ddot\theta=D^{-1}(u_{\mathrm{joint}}-C\dot\theta-G)$ تنها
$\theta_m$ ظاهر می‌شود؛ در $\theta^{(3)}$ تنها $\dot\theta_m$؛ و در $\theta^{(4)}$ جمله
$D^{-1}k\,\ddot\theta_m\ni D^{-1}\tfrac{k}{J_m}\tau$. پس درجه نسبی $4$ و ماتریس جداساز $D^{-1}k/J_m$، که
چون $D$ مثبت معین است هرگز تکین نمی‌شود. بازوی صلب بدون فنر، $D\ddot\theta=\tau-\ldots$، درجه
نسبی $2$ دارد.

\textbf{بسامد مُد موتور.} با ثابت نگه‌داشتن حلقه، $J_m\ddot\theta_m=-k\theta_m$، پس
$\omega=\sqrt{k/J_m}=\sqrt{100/0.1}=31.6$ رادیان بر ثانیه.

\textbf{نبود دینامیک صفر.} خطی‌سازی کامل $r\cdot n_y$ حالت را پوشش می‌دهد: $6\times1=6=n_x$ و
$4\times2=8=n_x$. هیچ بعدی باقی نمی‌ماند.

% ---------------------------------------------------------------------
\subsection{بازتولید}
\label{app:repro5}
نام توابع و فایل‌ها تنها در این جدول می‌آیند. همه اجراها در متلب و از ریشه مخزن انجام
می‌شوند؛ پوشه \en{Results/} در مخزن نیست.

\begin{latin}
\footnotesize
\setlength{\tabcolsep}{4pt}
\begin{tabular}{p{0.30\textwidth}p{0.30\textwidth}p{0.34\textwidth}}
    \toprule
    Result & Run & Output \\
    \midrule
    Spring-mass, pendulum, relative degree
        (\ref{tab:springmass}, \ref{tab:pendulum}, \ref{tab:reldeg}) &
        \texttt{ch5\_main} (studies \texttt{springmass pendulum relative\_degree}), summarized by \texttt{ch5\_report} &
        \texttt{Results/ch5\_<date>/} \\
    Admissibility at $x_0$ (\ref{tab:adm}) &
        \texttt{ch5\_ecbf\_admissible} & \texttt{Results/reruns/ch5\_admissibility.log} \\
    Pendulum ensemble (\ref{tab:pendens}) &
        \texttt{ch5\_ensemble} & \texttt{Results/reruns/ch5/} \\
    Moving $x_0$ (\ref{tab:moving}) &
        \texttt{ch5\_moving\_x0} & \texttt{Results/reruns/ch5\_moving\_x0/moving\_x0.log} \\
    Figures &
        \texttt{ch5\_doc\_figures} & \texttt{docs/figures/ch5\_*.png} \\
    Ensemble failures, dimensional bound &
        \texttt{ch5\_ensemble\_probe} (\emph{pending}) & \texttt{Results/reruns/ch5\_probe/} \\
    Joint-limit barrier with the height barrier &
        \texttt{ch5\_joint\_limit\_study} (\emph{pending}) & \texttt{Results/reruns/ch5\_joint\_limit/} \\
    Corollary 5.2 at the reference gait's post-impact state &
        \texttt{ch5\_walking\_preview} (\emph{pending}) & \texttt{Results/reruns/ch5\_walking\_preview/} \\
    Test suite (\ref{tab:tests}) &
        \texttt{ch5\_test\_all} & printed \\
    \bottomrule
\end{tabular}
\end{latin}

\end{document}
